\section{Source Map and Runtime Ownership}
Table~\ref{tab:source_map} identifies the principal implementation modules. The gateway is intentionally split between HTTP orchestration, scope policy, candidate selection, vector storage, and provider adapters. This avoids treating a successful route handler as proof that the retrieval or verification policy is correct. Each component has a narrower contract that can be tested independently.

\begin{table}[htbp]\centering\small
\caption{Code-grounded module map at the recorded source checkpoint.}\label{tab:source_map}
\begin{tabular}{p{4.7cm}p{7.8cm}}\toprule
Module & Responsibility and main output \\\midrule
\texttt{main.py}, \texttt{schemas.py}, \texttt{limits.py} & HTTP lifecycle; single-turn request validation; bounded body/concurrency; JSON or stream response. \\
\texttt{organization.py}, \texttt{identity.py} & Credential-to-principal lookup, permitted scope, provider profile, opaque identity. \\
\texttt{semantic\_cache.py} & Exact shortcut, vector lookup, top-$k$ decisions, final entry check, and scoped store. \\
\texttt{redisearch\_store.py}, \texttt{storage.py} & Redis hash/index access, atomic membership, TTL, LRU, byte and entry quotas. \\
\texttt{embeddings.py} & Lazy pinned MiniLM encoder or deterministic lexical hash backend. \\
\texttt{answer\_policy.py}, \texttt{verifier.py} & Candidate structural decisions and pinned NLI answer-assertion score. \\
\texttt{applicability.py}, \texttt{technical\_identity.py} & Optional bounded action, environment, and identity checks. \\
\texttt{providers/registry.py}, adapters & Intended route, breaker, first-token/idle deadlines, explicitly configured failover. \\
\texttt{singleflight.py}, \texttt{metrics.py} & Local duplicate JSON miss coalescing and operational counters/histograms. \\
\texttt{static/dashboard.html} & Optional operator/demo UI that constructs requests and displays actual outcomes. \\\bottomrule
\end{tabular}
\end{table}

\section{Startup, Readiness, and Local Versus Network Work}
The settings model reads environment values and an optional local environment file. On lifespan startup, the service creates text and binary Redis clients and checks connectivity. An organization directory can load hashed gateway-key digests and provider profiles. The gateway creates one embedder and semantic-cache coordinator for that process, builds profile-specific provider registries, and initializes rate limiting and metrics. If semantic mode uses MiniLM, a warmup encode can force the model to load. When the answer policy requires prewarming, failure to initialize the verifier prevents readiness. RedisSearch semantic mode probes the vector dimension, creates or validates the HNSW index, and fails if the existing vector schema is incompatible.

MiniLM and NLI inference execute locally in the gateway process and do not make a paid provider call. They can consume substantial CPU time and resident memory. Redis operations use a network protocol to the selected Redis server, even when it runs on the same machine. Only a miss or bypass invokes the configured upstream LLM adapter. With the development mock or controlled provider, that provider is local and deterministic. With an OpenAI-compatible or Gemini profile, the adapter makes an external HTTP call using a server-owned credential. The gateway cannot infer that such an adapter works with every endpoint merely from its construction; a bounded real-provider smoke remains a separate evidence gate.

Liveness is a minimal process-level status. Readiness includes selected dependencies and model/index state, while production validation also rejects unsafe default secrets and missing shared Redis configuration. Operational status can change after startup, so a startup check does not replace ongoing monitoring. Prewarming mainly moves a cold model load before readiness. It does not eliminate the computational cost of each new assertion or embedding.

\section{Request Schema and Authentication}
The supported request is single-turn. It supplies either a prompt of at most 16,000 characters or one user message with the same content limit. Extra unsupported fields are forbidden. Optional fields cover cache bypass, a private/workspace/none scope request, and a workspace identifier within its length bound. The body middleware also limits buffered bytes and active requests. These validation decisions make failure visible; silently accepting a full conversation history or unsupported generation settings would misrepresent the service's contract.

An application sends its gateway key in the \texttt{x-api-key} header. The directory compares its digest and resolves a principal rather than using the key text as a workspace identifier. That principal's organization, subject, application, memberships, permission set, and provider profile determine the allowed cache scope. The provider profile's secret comes from server configuration. An operator may rotate or revoke a gateway credential independently of the upstream provider credential. The handler rechecks the principal before returning a slow HIT and before provider admission. After generation, the recorded checkpoint checks authorization before cache storage: a revoked principal suppresses the write, but the completed JSON answer can still be returned. This is not a guarantee that revocation cancels an in-flight response.

As an illustrative API input, a trusted application can send
\begin{lstlisting}
POST /v1/chat
x-api-key: example-gateway-key
{"prompt":"How can I inspect the staging deployment?"}
\end{lstlisting}
The omitted scope fields mean automatic server-side selection; the example credential is fictitious. A developer-only override may request \texttt{cache\_scope=private}. An unauthorized explicit workspace request receives an authorization error rather than creating access through the API field. A normal dashboard request follows the automatic path. The inspected dashboard also exposes explicit controls in a separate developer mode and labels mock placeholders distinctly from cached or generated answers. Those UI observations refer to the committed dashboard at checkpoint \texttt{1cdb80a}; later UI changes require separate reconciliation.

\section{Embedding and Scoped Retrieval}
MiniLM converts each question into a 384-dimensional normalized FLOAT32 vector. An embedding is a numeric representation intended to put related text near one another. Cosine similarity is the dot product of normalized vectors; it ranks how close the question representations are. A high score is compatible with a changed action or environment and is never, by itself, an answer-approval rule. The test-only hash backend instead counts normalized tokens into a deterministic vector; it is fast and dependency-light but lexical rather than a substitute for the measured semantic model.

In semantic RedisSearch mode, the index has a TAG field for the opaque scope and an HNSW vector field for cosine distance. The query requests a bounded top-$k$ set, with limited overfetch to compensate when a returned hash has expired. Each search result is matched against its current Redis hash for scope, question, and answer identity before entering the candidate list. A process-local matrix is the alternative backend. It still reads Redis as the authoritative store but keeps a local vector snapshot and locks around rebuilds. A NumPy dot-product scan of that matrix is linear in entry count, not a proven sub-linear index. The shared RedisSearch index is more suitable for multiple workers, but its approximate retrieval recall and large-scale capacity have not been established for this project.

The exact shortcut computes a question digest and reads its scope-local Redis hash before embedding. This is an optimization for identical text; it is not the main mechanism for paraphrase reuse. It rereads after updating recency so an expired or replaced entry is not served from an earlier local view. In semantic mode, candidates are evaluated in rank order. A lower-ranked candidate can win if higher-ranked candidates are declined, subject to the limit, deadline, and conflict rule.

\section{Answer Suitability and Selection}
For each non-exact candidate the answer policy receives the incoming question, stored question, stored answer, and retrieval similarity. Its first checks reject malformed or too-distant entries and changed numbers, polarity, named works, and answer type. Optional technical and applicability checks examine bounded syntax such as requested action, HTTP method, target environment, version, permission, or named workspace when those facets are supported and grounded in the recorded answer. An unsupported or ambiguous phrasing can abstain. These checks deliberately trade some valid reuse for a more explicit reason to decline a changed task; they do not understand arbitrary technical requests.

An unresolved candidate is sent to a pinned local NLI cross-encoder as a pair of assertions: the stored answer is asserted for the stored question and then for the incoming question. The model's entailment probability is a learned signal about compatibility of those statements. It does not verify that the stored answer was factually true, current, complete, or authorized for a new audience. Paired inputs are tokenized without silent truncation; if they exceed the model window, the policy returns VERIFY. Unavailable inference also produces VERIFY. The policy artifact supplies distinct accept and reject thresholds, leaving an uncertainty interval. A candidate action of VERIFY means abstain and use the configured provider path if no other candidate is accepted; it does not call a human reviewer.

The selector records a trace of each considered rank, similarity, action, reason, score status, and elapsed time. A configured verification limit and millisecond budget cap work on the critical path. If plausible lower candidates remain unassessed, or accepted candidates disagree on answer text, the entire request becomes a MISS. After choosing an accepted candidate, the coordinator reads the original hash again and compares the question and answer. A disappeared or changed entry is not served. This last check addresses a practical race between vector search, model inference, expiry, and concurrent writes.

\section{Provider Completion, Retention, and Errors}
The registry resolves an intended model route from the principal's provider profile. Optional complexity routing can choose a cheaper or stronger model, but that route is an engineering heuristic and its answer quality has not been validated. Circuit-breaker tickets track transient failures; the registry can try only providers that were explicitly configured. It never invents a mock fallback for a real provider failure. First-token, idle-stream, provider, lookup, and overall request deadlines protect different stages. An exceeded overall deadline can prevent provider dispatch entirely.

JSON generation assembles content events and requires one normal completion with a stop reason. Empty, oversized, interrupted, or non-normal completions return errors and are not cached. Streaming emits content and metadata events and only emits a final completion marker after success. An interrupted stream must not become a reusable partial answer. A completed response is classified for storage. A request or response marked no-store is not written. A private-only response from a workspace request uses a private scope. A successful fallback through a different provider is not stored under the intended provider's compatibility identity. A storage failure is reported while retaining the already completed generated answer.

The cache TTL, scope LRU, and organization count/byte quotas constrain residence and memory usage in different ways. TTL is time since insertion; LRU chooses a victim under pressure. Occupancy reports are useful for operations but do not measure answer quality. Rate limiting applies before lookup, while the daily admitted-miss count is reached only after a miss or bypass. Because it counts requests rather than billed tokens, it cannot enforce a monetary spending cap.

\section{Four Request Walkthroughs}
\textbf{First request (illustrative).} Alice's integrated application sends an eligible technical question with no scope override. Her principal resolves to Organization A's authorized workspace and its provider profile. No current exact or semantic candidate is approved, so the request is admitted for a provider call. If the configured provider completes normally and its answer is shareable, the gateway stores question, answer, vector, metadata, and TTL under the compatible workspace identity. The JSON response reports MISS and the provider name. If a cache write fails, Alice still receives the complete generated answer with a failure storage status.

\textbf{Colleague's paraphrase (illustrative).} Carol has membership in the same workspace and compatible profile. Her differently worded request misses the internal exact shortcut. MiniLM produces a vector, RedisSearch returns Alice's stored entry, and the policy checks the particular answer against Carol's task. If it accepts and the final hash is unchanged, Carol receives a HIT without provider generation. The saved isolated release probe recorded such a seeded decorator-answer hit, but that selected example does not establish general answer accuracy.

\textbf{Changed task (illustrative).} Alice's stored answer explains how to preview a Terraform plan for staging. Carol asks to apply a production change. The questions can be near each other in vector space, yet the action and environment differ. An enabled applicability guard can reject the candidate before model scoring; if a guard is off or cannot parse the request, the NLI path may still accept or abstain. A declined candidate yields a provider miss, which is desirable when reuse would be unsuitable. The controlled evaluations in Chapter~6 document both prevented and missed cases.

\textbf{Private or no-store request (illustrative).} Alice asks about account-specific information. A private explicit scope or the server-side eligibility decision routes storage into Alice's subject/application namespace. Carol's workspace lookup cannot read it, and Organization B has an entirely different scope identity. A no-store decision bypasses cache reuse and suppresses storage after generation. This is a bounded sharing policy, not a guarantee that a heuristic identifies every sensitive string.
